% Part: proof-theory
% Chapter: sequent-calculus
% Section: proof-examples

\documentclass[../../../include/open-logic-section]{subfiles}

\begin{document}

\olfileid[id]{pt}{seq}{ex}

\olsection{Contoh Pembuktian}

Untuk memberikan !!{proof}s atas sekuen, biasanya lebih mudah bekerja
dari sekuen akhir ke arah atas: pilih !!a{formula} di dalamnya sebagai
!!{formula} aktif, tuliskan premis-premis aturan yang bersesuaian di atas
sekuen itu, lalu ulangi proses tersebut sampai menghasilkan aksioma.
Sebagai contoh, andaikan kita hendak membuktikan sekuen
\[(!C \land !D) \lif !E \Sequent !C \lif (!D \lif !E)\]
Tinjau $!C \lif (!D \lif !E)$: formula ini berbentuk $!A \lif !B$ dan
muncul di sebelah kanan sekuen. Mencocokkan seluruh sekuen dengan aturan
$\RightR{\lif}$ dari \Log{G1c},
\[\Axiom$ !A, \Gamma \fCenter \Delta, !B$
\RightLabel{\RightR{\lif}}
\UnaryInf$ \Gamma \fCenter \Delta, !A \lif !B $
\DisplayProof\]
menyarankan agar kita mengambil $\Gamma = \{(!C \land !D) \lif !E\}$,
$\Delta = \emptyset$, $!A \ident !C$, dan $!B \ident !D \lif !E$.
Jadi, inferensi terakhir suatu bukti dapat berupa
\[\Axiom$ !C, (!C \land !D) \lif !E \fCenter !D \lif !E$
\RightLabel{\RightR{\lif}}
\UnaryInf$(!C \land !D) \lif !E \fCenter !C \lif (!D \lif !E)$
\DisplayProof\]
Jika gagasan ini kita ulangi dengan memusatkan perhatian pada $!D \lif
!E$ sebagai formula aktif, kita memperoleh
\[\Axiom$ !D, !C, (!C \land !D) \lif !E \fCenter !E$
\RightLabel{\RightR{\lif}}
\UnaryInf$!C, (!C \land !D) \lif !E \fCenter !D \lif !E$
\RightLabel{\RightR{\lif}}
\UnaryInf$(!C \land !D) \lif !E \fCenter !C \lif (!D \lif !E)$
\DisplayProof\]
Sekarang kita harus memusatkan perhatian pada $(!C \land !D) \lif !E$
dan memakai aturan $\LeftR{\lif}$ dari \Log{G1c},
\[\Axiom$ \Gamma \fCenter \Delta, !A$
\Axiom$ !B, \Gamma \fCenter \Delta$
\RightLabel{\LeftR{\lif}}
\BinaryInf$ !A \lif !B, \Gamma \fCenter \Delta$
\DisplayProof\]
Hal ini menghasilkan
\[\Axiom$!D, !C \fCenter !E, !C \land !D$
\Axiom$!D, !C, !E \fCenter !E$
\RightLabel{\LeftR{\lif}}
\BinaryInf$ !D, !C, (!C \land !D) \lif !E \fCenter !E$
\RightLabel{\RightR{\lif}}
\UnaryInf$!C, (!C \land !D) \lif !E \fCenter !D \lif !E$
\RightLabel{\RightR{\lif}}
\UnaryInf$(!C \land !D) \lif !E \fCenter !C \lif (!D \lif !E)$
\DisplayProof\]
Dengan menjadikan !!{formula} $!C \land !D$ pada sekuen paling kiri di
baris teratas sebagai formula prinsipal dan menggunakan aturan
$\RightR{\land}$, kita memperoleh
\[
\Axiom$!D, !C \fCenter !E, !C$
\Axiom$!D, !C \fCenter !E, !D$
\RightLabel{\RightR{\land}}
\BinaryInf$!D, !C \fCenter !E, !C \land !D$
\Axiom$!D, !C, !E \fCenter !E$
\RightLabel{\LeftR{\lif}}
\BinaryInf$ !D, !C, (!C \land !D) \lif !E \fCenter !E$
\RightLabel{\RightR{\lif}}
\UnaryInf$!C, (!C \land !D) \lif !E \fCenter !D \lif !E$
\RightLabel{\RightR{\lif}}
\UnaryInf$(!C \land !D) \lif !E \fCenter !C \lif (!D \lif !E)$
\DisplayProof
\]
Sejauh ini semuanya berjalan baik. Perhatikan bahwa segala sesuatu yang
kita lakukan juga berlaku dalam \Log{G3c}, karena aturan yang sejauh ini
kita gunakan sama dalam \Log{G1c} dan~\Log{G3c}. Kita telah sampai pada
sebuah fragmen !!a{proof} yang semua sekuen teratasnya memuat
!!a{formula} yang sama di sebelah kiri dan kanan. Namun, sekuen-sekuen
tersebut belum benar-benar merupakan aksioma.

Dalam \Log{G1c}, kita harus memberikan !!{proof}s atas sekuen-sekuen
teratas dari aksioma berbentuk $!A \Sequent !A$. Hal ini mudah dilakukan
dengan aturan pelemahan. Sebagai contoh, kita dapat !!{prove} sekuen teratas
paling kiri $!D, !C \Sequent !E, !C$ sebagai berikut:
\[
\Axiom$!C \fCenter !C$
\RightLabel{\LeftR{\Weakening}}
\UnaryInf$!D, !C \fCenter !C$
\RightLabel{\RightR{\Weakening}}
\UnaryInf$!D, !C \fCenter !E, !C$
\DisplayProof
\]
Secara keseluruhan, kita mempunyai !!a{proof} dalam \Log{G1c} yang
berbentuk sebagai berikut:
\[
\Axiom$!C \fCenter !C$
\RightLabel{\LeftR{\Weakening}}
\UnaryInf$!D, !C \fCenter !C$
\RightLabel{\RightR{\Weakening}}
\UnaryInf$!D, !C \fCenter !E, !C$
\Axiom$!D \fCenter !D$
\RightLabel{\LeftR{\Weakening}}
\UnaryInf$!D, !C \fCenter !D$
\RightLabel{\RightR{\Weakening}}
\UnaryInf$!D, !C \fCenter !E, !D$
\RightLabel{\RightR{\land}}
\BinaryInf$!D, !C \fCenter !E, !C \land !D$
\Axiom$!E \fCenter !E$
\RightLabel{\LeftR{\Weakening}}
\UnaryInf$!C, !E \fCenter !E$
\RightLabel{\LeftR{\Weakening}}
\UnaryInf$!D, !C, !E \fCenter !E$
\RightLabel{\LeftR{\lif}}
\BinaryInf$ !D, !C, (!C \land !D) \lif !E \fCenter !E$
\RightLabel{\RightR{\lif}}
\UnaryInf$!C, (!C \land !D) \lif !E \fCenter !D \lif !E$
\RightLabel{\RightR{\lif}}
\UnaryInf$(!C \land !D) \lif !E \fCenter !C \lif (!D \lif !E)$
\DisplayProof
\]
Struktur bukti ini (khususnya ketinggiannya) tidak bergantung pada apa
!!{formula}s $!C$, $!D$, dan~$!E$ itu. Kita dapat mengganti simbol-simbol
tersebut dalam !!{proof} di atas dengan !!{formula}s apa pun dan hasilnya
tetap merupakan !!a{proof} dalam~\Log{G1c}. Bukti dalam \Log{G1c} itu
bersifat \emph{skematis}.

Dalam \Log{G3c}, aksioma adalah sekuen berbentuk $!A, \Gamma \Sequent
\Delta, !A$ dengan syarat $!A$ atomik. Jadi, meskipun $!D, !C \Sequent
!E, !C$ tampak seperti aksioma \Log{G3c} ($\Gamma = \{!D\}$; $\Delta =
\{!E\}$), sekuen tersebut \emph{benar-benar} merupakan aksioma hanya jika
$!C$ memang atomik. Fragmen bukti kita menjadi bukti lengkap dalam
\Log{G3c} hanya jika $!C$, $!D$, dan $!E$ merupakan !!{formula}s atomik.

Meskipun, secara tegas, kita belum menunjukkan bahwa
\[\Log{G3c} \Proves (!C \land !D) \lif !E \fCenter !C \lif (!D \lif
!E),\] dalam praktiknya inilah semua yang perlu (dan dapat) kita lakukan.
Alasannya ialah bahwa setiap sekuen berbentuk $!A, \Gamma \Sequent
\Delta, !A$ dapat dibuktikan dalam \Log{G3c}, bahkan jika $!A$ sendiri
tidak atomik. Kita dapat menunjukkan hal ini dengan induksi pada
kedalaman~$\depth{!A}$ dari~$!A$.

\begin{defn}\ollabel{defn:depth}
\emph{Kedalaman}~$\depth{!A}$ dari !!a{formula} didefinisikan sebagai berikut:
\begin{enumerate}
  \item Jika $!A$ atomik, $\depth{!A} = 0$.
  \item Jika $!A \ident \lnot !B$, $!A \ident \lforall[x][!B(x)]$,
  atau $!A \ident \lexists[x][!B(x)]$, maka $\depth{!A} = \depth{!B} + 1$.
  \item Jika $!A \ident (!B \land !C)$, $!A \ident (!B \lor !C)$, atau
  $!A \ident (!B \lif !C)$, maka $\depth{!A} =
  \max(\depth{!B},\depth{!C}) + 1$.
\end{enumerate}
\end{defn}

Tidak sulit membuktikan bahwa $\depth{A(x)} = \depth{A(t)}$ untuk setiap
term~$t$. Fakta yang penting bagi kita ialah bahwa, dalam setiap aturan,
kedalaman !!{formula} prinsipal selalu lebih besar daripada kedalaman
!!{formula}s aktif. Sebagai contoh, kita mempunyai:
\begin{align*}
\depth{P(x)} = \depth{P(a)} & = 0,\\
\depth{\lforall[x][P(x)]} & = 1,\\
\depth{(\lforall[x][P(x)] \lif P(a))} & = \max(1,0) + 1 = 2.
\end{align*}
Kita dapat menggunakan fakta ini untuk membuktikan beberapa hasil dengan
induksi pada~$\depth{!A}$, misalnya hasil berikut.

\begin{prop}\ollabel{prop:G1c-axioms}
Untuk setiap~$!A$, terdapat !!a{proof} dalam $\Log{G1c}$ atas $!A
\Sequent !A$ yang semua aksiomanya atomik.
\end{prop}

\begin{proof}
  Dengan induksi pada~$\depth{!A}$. Jika $\depth{!A} = 0$, maka $!A$
  atomik dan $!A \Sequent !A$ sendiri sudah merupakan !!a{proof}
  dalam~\Log{G1c} dengan bentuk yang diperlukan. Jika $\depth{!A} > 0$,
  $!A$ dibangun dari !!{formula}s yang berkedalaman lebih rendah, misalnya
  $!A \ident \lnot !B$, $!A \ident (!B \land !C)$, atau $!A \ident
  \lforall[x][!B(x)]$. Hipotesis induksinya ialah bahwa, untuk semua $!D$
  dengan $\depth{!D} < \depth{!A}$, $\Log{G1c} \Proves !D \Sequent !D$
  dari aksioma-aksioma atomik.

  Jika $!A \ident \lnot !B$, maka $\depth{!B} < \depth{!A}$ dan, menurut
  hipotesis induksi, terdapat !!a{proof}~$\pi_1$ dalam \Log{G1c} dari
  aksioma atomik atas $!B \Sequent !B$. Kita dapat memperluasnya menjadi
  !!a{proof} atas $\lnot !B \Sequent \lnot !B$ sebagai berikut:
  \[
  \AxiomC{}
  \RightLabel{$\pi_1$}
  \Deduce$!B \fCenter !B$
  \RightLabel{\LeftR{\lnot}}
  \UnaryInf$\lnot !B, !B \fCenter $
  \RightLabel{\RightR{\lnot}}
  \UnaryInf$\lnot !B \fCenter \lnot !B$
  \DisplayProof
\]

  Jika $!A \ident (!B \land !C)$, maka $\depth{!B} < \depth{!A}$ dan
  $\depth{!C} < \depth{!A}$. Menurut hipotesis induksi, terdapat
  !!{proof}s~$\pi_1$ dan~$\pi_2$ dalam \Log{G1c} dari aksioma atomik
  masing-masing atas $!B \Sequent !B$ dan $!C \Sequent !C$. Kita dapat
  memperluasnya menjadi !!a{proof} atas $!B \land !C \Sequent !B \land
  !C$ sebagai berikut:
  \[
  \AxiomC{}
  \RightLabel{$\pi_1$}
  \Deduce$!B \fCenter !B$
  \RightLabel{\LeftR{\Weakening}}
  \UnaryInf$!B, !C, \fCenter !B$
  \RightLabel{\LeftR{\land}}
  \UnaryInf$!B \land !C, !C\fCenter !B$
  \RightLabel{\LeftR{\land}}
  \UnaryInf$!B \land !C, !B \land !C \fCenter !B$
  \AxiomC{}
  \RightLabel{$\pi_2$}
  \Deduce$!C \fCenter !C$
  \RightLabel{\LeftR{\Weakening}}
  \UnaryInf$!B, !C, \fCenter !C$
  \RightLabel{\LeftR{\land}}
  \UnaryInf$!B \land !C, !C \fCenter !C$
  \RightLabel{\LeftR{\land}}
  \UnaryInf$!B \land !C, !B \land !C \fCenter !C$
  \RightLabel{\RightR{\land}}
  \BinaryInf$!B \land !C, !B \land !C \fCenter !B \land !C$
  \RightLabel{\LeftR{\Contraction}}
  \UnaryInf$!B \land !C \fCenter !B \land !C$
  \DisplayProof
  \]

  Sekarang andaikan $!A \ident \lforall[x][B(x)]$. Ambil
  !!a{constant}~$c$ yang tidak muncul dalam $\lforall[x][B(x)]$. Maka
  $\depth{B(c)} = \depth{!B(x)} < \depth{!A}$ dan, menurut hipotesis
  induksi, terdapat !!a{proof}~$\pi_1$ dalam \Log{G1c} dari aksioma atomik
  atas $!B(c) \Sequent !B(c)$. Kita dapat memperluasnya menjadi
  !!a{proof} atas $\lforall[x][B(x)] \Sequent \lforall[x][B(x)]$
  sebagai berikut:
  \[
  \AxiomC{}
  \RightLabel{$\pi_1$}
  \Deduce$!B(c) \fCenter !B(c)$
  \RightLabel{\LeftR{\lforall}}
  \UnaryInf$\lforall[x][!B(x)] \fCenter !B(c)$
  \RightLabel{\RightR{\lforall}}
  \UnaryInf$\lforall[x][!B(x)] \fCenter \lforall[x][!B(x)]$
  \DisplayProof
\]

  Kasus-kasus yang tersisa diserahkan sebagai latihan.
\end{proof}

\begin{prob}
  Lengkapi bukti \olref{prop:G1c-axioms} dengan mengerjakan kasus-kasus
  ketika $!A \ident (!B \lor !C)$, $!A \ident (!B \lif !C)$, dan
  $!A \ident \lexists[x][B(x)]$.
\end{prob}

Untuk keperluan kita, dalam kasus $!A \ident \lnot !B$ kita juga dapat
menggunakan !!{proof}
\[
\AxiomC{}
\RightLabel{$\pi_1$}
\Deduce$!B \fCenter !B$
\RightLabel{\RightR{\lnot}}
\UnaryInf$\fCenter !B, \lnot !B$
\RightLabel{\LeftR{\lnot}}
\UnaryInf$\lnot !B \fCenter \lnot !B$
\DisplayProof
\]
Namun, ada sistem sekuen yang tidak memungkinkan cara ini. Kalkulus sekuen
intuisionistik \Log{G1i} serupa dengan \Log{G1c}, kecuali bahwa semua
sekuennya dibatasi agar memuat paling banyak satu !!a{formula} di sebelah
kanan. Jika $\Delta = \emptyset$, !!{proof} pertama juga merupakan
!!a{proof} dalam~\Log{G1i}, tetapi yang kedua bukan, sebab salah satu
sekuennya memuat $!B$ sekaligus~$\lnot !B$ di sebelah kanan.

Perhatikan bahwa bahkan dalam \Log{G1c} kita tidak dapat menggunakan urutan
aturan yang berbeda dalam kasus $!A \ident \lforall[x][!B(x)]$. Yang
berikut \emph{bukan} !!a{proof}:
  \[
  \AxiomC{}
  \RightLabel{$\pi_1$}
  \Deduce$!B(c) \fCenter !B(c)$
  \RightLabel{\RightR{\lforall}}
  \UnaryInf$!B(c) \fCenter \lforall[x][!B(x)]$
  \RightLabel{\LeftR{\lforall}}
  \UnaryInf$\lforall[x][!B(x)] \fCenter \lforall[x][!B(x)]$
  \DisplayProof
\]
karena syarat variabel eigen pada~$\RightR{\lforall}$ dilanggar.

\begin{prop}\ollabel{prop:G3c-axioms}
Setiap sekuen berbentuk $!A, \Gamma \Sequent \Delta, !A$ (dengan $!A$
yang \emph{tidak} harus atomik) dapat dibuktikan dalam~\Log{G3c}.
\end{prop}

\begin{proof}
Buktinya sangat mirip dengan bukti \olref{prop:G1c-axioms}, tetapi kita
harus berhati-hati dengan hipotesis induksi. Kasus $n = \depth{!A} = 0$
kembali mudah. Jika $n=0$, $!A$ atomik, dan $!A, \Gamma \Sequent \Delta,
!A$ merupakan aksioma~\Log{G3c}.

Sekarang andaikan $n>0$. Kita kembali membedakan beberapa kasus. Hipotesis
induksinya ialah: untuk setiap $!D$ dengan $\depth{!D} < n$, serta setiap
$\Pi$ dan $\Lambda$, $\Log{G3c} \Proves !D, \Pi \Sequent \Lambda, !D$.
Berikut kasus untuk $!A \ident (!B \land !C)$. Kita menggunakan hipotesis
induksi dua kali: pertama untuk $!D = !B$, $\Pi = !C, \Gamma$, dan
$\Lambda = \Delta$ sehingga diperoleh !!a{proof}~$\pi_1$ atas $!B, !C,
\Gamma \Sequent \Delta, !B$; kedua untuk $!D = !C$, $\Pi = !B, \Gamma$,
dan $\Lambda = \Delta$ sehingga diperoleh !!a{proof}~$\pi_2$ atas $!B,
!C, \Gamma \Sequent \Delta, !C$. Berikut !!{proof} atas~$!B \land !C,
\Gamma \Sequent \Delta, !B \land !C$:
\[
\AxiomC{}
\RightLabel{$\pi_1$}
\Deduce$!B, !C, \Gamma \fCenter \Delta, !B$
\RightLabel{\LeftR{\land}}
\UnaryInf$!B \land !C, \Gamma \fCenter \Delta, !B$
\AxiomC{}
\RightLabel{$\pi_2$}
\Deduce$!B, !C, \Gamma \fCenter \Delta, !C$
\RightLabel{\LeftR{\land}}
\UnaryInf$!B \land !C, \Gamma \fCenter \Delta, !C$
\RightLabel{\RightR{\land}}
\BinaryInf$!B \land !C, \Gamma \fCenter \Delta, !B \land !C$
\DisplayProof
\]

Untuk $!A \ident \lforall[x][!B(x)]$, pilih !!a{constant} yang tidak
muncul dalam~$\lforall[x][!B(x)]$, $\Gamma$, ataupun~$\Delta$. Terapkan
hipotesis induksi dengan $!D = !B(c)$, $\Pi = \lforall[x][!B(x)], \Gamma$,
dan $\Lambda = \Delta$. Hal ini menghasilkan !!a{proof}~$\pi_1$ atas
$!B(c), \lforall[x][!B(x)], \Gamma \Sequent \Delta, !B(c)$, yang dapat
kita lanjutkan sebagai berikut:
\[
\AxiomC{}
\RightLabel{$\pi_1$}
\Deduce$!B(c), \lforall[x][!B(x)], \Gamma \fCenter \Delta, !B(c)$
\RightLabel{\LeftR{\lforall}}
\UnaryInf$\lforall[x][!B(x)], \Gamma \fCenter \Delta, !B(c)$
\RightLabel{\RightR{\lforall}}
\UnaryInf$\lforall[x][!B(x)], \Gamma \fCenter \Delta, \lforall[x][!B(x)]$
\DisplayProof
\]
Syarat variabel eigen pada $\RightR{\lforall}$ terpenuhi berkat cara kita
memilih~$c$.
\end{proof}

\begin{prob}
  Lanjutkan bukti \olref{prop:G3c-axioms} dengan mengerjakan kasus-kasus
  ketika $!A \ident (!B \lif !C)$ dan $!A \ident \lexists[x][B(x)]$.
\end{prob}

\end{document}
